\section{Section-ratio fields and horizontal regularity}
\label{g:cf:fields}

The contraction in Theorem~\ref{g:cf:contraction-theorem} uses a
maximal linear system. There is an
intrinsic description of its base field: it consists of the rational
functions whose zeros and poles have degree zero against the supporting
movable class. This identification lets us compare two ways of producing
the new generic section. Affine interpolation uses the precise relation
between three twists. A second argument uses finite horizontal support
and the ordering of nonnegative integer vectors. We first separate the
regularity and uniqueness conditions that both arguments must preserve.

\subsection{Regularity on the original variety and on a resolution}
\label{g:cf:conditions}

Let $X$ be smooth connected projective, let an algebraic torus $T$ act on
$X$, and let $L$ be a $T$-linearized line bundle. Fix a relatively
algebraically closed subfield
$\mathbb C\subseteq K\subseteq\mathbb C(X)^T$. A divisorial valuation is
\emph{horizontal over $K$} when it is zero on $K^*$; on a proper model
mapping to a projective model of $K$, this means that its center dominates
the base. All orders of sections on higher models refer to pullbacks of
their line bundles.

Suppose $s$ is an invariant rational section of $aL$, where $a>0$, whose
pullback is regular on the complete generic fiber of a resolution of the
map defined by $K$. This is a condition on every horizontal divisorial
valuation, including exceptional ones. It is unchanged by further
resolution, because a regular section pulls back regularly; the converse
follows by pushforward and normality. The uniqueness condition below,
in contrast, tests the other rational sections only at prime divisors of
$X$ itself.

\begin{lemma}[Two equivalent uniqueness conditions]
\label{g:cf:uniqueness}
With these data, the following statements are equivalent.
\begin{enumerate}
\item For every integer $j\geq0$ and every nonzero invariant rational
section $u$ of $jaL$ with no poles at the $K$-horizontal primes of $X$,
one has $u/s^j\in K$.
\item For every integer $b\geq0$, the ratio of any two nonzero invariant
rational sections of $bL$ having no poles at the $K$-horizontal primes of
$X$ belongs to $K$.
\end{enumerate}
Both statements include degree zero.
\end{lemma}

\begin{proof}
The second assertion implies the first by comparing $u$ with $s^j$;
the regularity hypothesis on $s$ guarantees that this comparison is
permitted. Conversely, let $u_1,u_2$ have degree $b$. The first assertion
applied to $u_i^a$ gives $u_i^a/s^b\in K$. Hence
$(u_1/u_2)^a\in K$, and relative algebraic closedness puts $u_1/u_2$ in
$K$. This argument also applies when $b=0$.
\end{proof}

These equivalent formulations do not identify regularity at original
primes with regularity at all upstairs primes. The latter is the stronger
regularity requirement on the chosen section $s$. In particular, a
uniqueness assertion that tests only sections already regular at every
upstairs horizontal prime is not, without another argument, the
uniqueness assertion in Lemma~\ref{g:cf:uniqueness}.

For $L=-K_X$, these data imply the adjoint rank hypotheses of
Theorem~\ref{g:cf:contraction-theorem}. Indeed, on a smooth resolution
$\pi:Z\to X$, $g:Z\to Y$, put $R=K_{Z/X}$. Dividing an invariant
section of $R+ja\pi^*L$ over the generic point of $Y$ by the Jacobian
section gives a rational section of $jaL$ without poles at horizontal
primes of $X$: their strict transforms have coefficient zero in $R$.
Lemma~\ref{g:cf:uniqueness} makes it a $K$-multiple of $s^j$.
Conversely $s^j s_R$ is present. Thus the invariant rank is one for
every $j\geq0$. This implication uses the original-prime test; adjoint
rank alone does not license replacing it by a test on arbitrary rational
sections with exceptional poles.

\subsection{The null-prime field is the maximal-system field}

\begin{lemma}\label{g:cf:null-field}
Let $Y$ be a smooth connected projective variety of dimension $d>0$ and
let $\alpha\ne0$ be a movable numerical curve class. Call a prime divisor
$P$ \emph{null} if $\alpha\cdot P=0$, and set
\[
 K_\alpha=\{0\}\cup\{u\in\mathbb C(Y)^*:
          \operatorname{Supp}(\operatorname{div}u)
          \text{ consists of null primes}\}.
\]
Then $K_\alpha$ is relatively algebraically closed in $\mathbb C(Y)$
and has transcendence degree less than $d$. It is the field of the Stein
base of any linear system of degree zero against $\alpha$ having maximal
image dimension among all such systems. Every prime of $Y$ vertical over
$K_\alpha$ is null. If $\operatorname{Pic}(Y)$ is finitely generated,
only finitely many null primes are horizontal over $K_\alpha$.
\end{lemma}

\begin{proof}
Products and inverses preserve the support condition. The poles of a
nonzero sum remain on null primes. Since a principal divisor has degree
zero, its effective zero divisor then has degree zero as well; movability
makes each of its components null. Thus $K_\alpha$ is a field.

Equivalently, every valuation at a non-null prime is trivial on
$K_\alpha$. Such a valuation remains trivial on an element algebraic over
that field. An element of positive valuation cannot satisfy a polynomial
with nonzero constant term over the field, since that term would uniquely
have minimal valuation. Applying the same argument to the inverse excludes
negative valuation. This proves relative algebraic closedness.

If $u_1,\ldots,u_d\in K_\alpha$ were algebraically independent, the
sum $P$ of their polar divisors would be null. The monomials of total
degree at most $m$ in the $u_i$ are linearly independent sections of
$mP$, so $P$ would be big. A nonzero functional nonnegative on the
pseudoeffective cone is strictly positive on its interior. Thus
$\alpha\cdot P>0$, a contradiction.

Every ratio in a zero-degree system belongs to $K_\alpha$, since the
zero divisors of its sections are effective of degree zero. Conversely,
if $u\in K_\alpha^*$ and $P=(\operatorname{div}u)^-$, then $1$ and
$u$ are sections of $P$ and $\alpha\cdot P=0$. Choose a maximal-image
zero-degree system and write $K_0$ for the field of its Stein base.
The product-system argument of Lemma~\ref{g:cf:maximal-field} puts
all zero-degree section ratios in $K_0$, hence $K_\alpha\subseteq K_0$.
The generating ratios of the chosen system lie in $K_\alpha$; its
relative algebraic closedness gives the reverse inclusion. This proves
the claimed equality, rather than just a comparison of dimensions.

A non-null prime is horizontal over $K_\alpha$ by definition. Finally,
if there were infinitely many horizontal null primes, finite generation
of $\operatorname{Pic}(Y)$ would give a nonzero integral relation among
finitely many of their classes (multiply to remove torsion). The resulting
nonzero divisor would be principal. Its defining function would lie in
$K_\alpha$ but have nonzero order at one of these horizontal primes,
which is impossible.
\end{proof}

The hypotheses for finiteness hold on every smooth base dominated by the
variety in Theorem~\ref{g:cf:contraction-theorem}: pullback and birational
invariance give $H^0(Y,\Omega_Y^1)=0$, hence $\operatorname{Pic}^0(Y)=0$.
Such subfields are finitely generated: choose a finite transcendence basis;
the relative algebraic closure of its rational function field inside the
finitely generated field $\mathbb C(Y)$ is finite. Consequently the field
has a smooth projective model. The same argument will apply to the fields
constructed below.

\subsection{Minima on all components and on original components}
\label{g:cf:minima}

We now retain the fixed-section data of
Section~\ref{g:cf:conditions}, with $L=-K_X$ smoothly semipositive and
$H^0(X,\Omega_X^q)=0$ for $q>0$. Assume also either of the equivalent
uniqueness conditions in Lemma~\ref{g:cf:uniqueness}. Let $Y$ be a smooth projective model
of $K$. Prepare a normal graph
\[
 X\xleftarrow{p}W\xrightarrow{f}Y
\]
whose fibers have dimension at most $\dim X-\dim Y$.
This preparation can be made without toroidalization. On the flat locus
of an initial graph, the fibers define a map to the projective Hilbert
scheme of $X$. Resolve that map on the base and pull back the universal
family. Its fixed Hilbert polynomial bounds the dimension of every fiber.
The main graph transform lies in this family, and its normalization is
finite over the transform. This gives the asserted bound. The graph and
its normalization are equivariant; take an equivariant
smooth resolution $\rho:Z\to W$ afterwards
\cite{RaynaudGruson1971}, and write $\pi=p\rho$ and $g_Z=f\rho$ for
its composite maps to $X$ and $Y$. In particular every vertical
prime of $W$ dominates a prime of $Y$.

For a prime $P\subset Y$, let $Q$ range over the components above it on
$W$, and put $e_Q=\operatorname{ord}_Q f^*P$. Define
\[
 d_P=\min_{Q\mapsto P}\frac{\operatorname{ord}_Q(s)}{e_Q},\qquad
 d_P^*=\min_{\substack{Q\mapsto P\\Q\text{ nonexceptional over }X}}
                       \frac{\operatorname{ord}_Q(s)}{e_Q}.
\]
A prime $P$ is \emph{omitted} if all components above it are exceptional
over $X$; set $d_P^*=d_P$ in this case. There are finitely many omitted
primes. The divisors
\[
 D=\sum_Pd_PP,\qquad D^*=\sum_Pd_P^*P
\]
have finite support, $D^*\geq D$, and
\[
 \operatorname{div}_W(s)-f^*D\geq0.
\]
The last divisor is rational Cartier, so it pulls back effectively to
$Z$. The first minimum controls the metric because some component has
zero normalized section order. The second controls sections on $X$
because a pole on an exceptional divisor disappears on pushdown.

\begin{lemma}\label{g:cf:finite-volume-base}
The rational divisor $D$ has a semipositive metric with locally bounded
weights. There is an effective integral divisor $P_0$, supported on the
omitted primes, and a psh metric weight $\phi$ on $-K_Y+P_0$ such that
$|s_{P_0}|^2e^{-\phi}$ is locally integrable.
\end{lemma}

\begin{proof}
If $Y$ is a point, then $D=0$; take $P_0=0$ and positive Hermitian
norms on the two trivial lines. Their weights are bounded and their
adjoint measure is finite. We may therefore assume $\dim Y>0$.
Apply Lemmas~\ref{g:moments} and~\ref{g:transfer} with original
variety $X$, allowed-pole divisor zero, $A=L=-K_X$, and both metrics equal
to the compact-averaged smooth metric on $L$. The section is $s$ of degree
$a$. It has no horizontal poles on the normal equidimensional graph.
The rank computation following Lemma~\ref{g:cf:uniqueness} verifies the
smooth invariant adjoint rank condition for every $j\geq0$. Thus the
two required psh functions are supplied by Lemma~\ref{g:moments}.

The full-component normalization in Lemma~\ref{g:transfer} is exactly
$D$; a denominator-clearing multiple of $D$ receives its bounded metric,
which scales back to the rational divisor. Since the allowed-pole divisor
on $X$ is zero, the invisible primes in that lemma are precisely our
omitted primes. Its correction is therefore an effective integral $P_0$
with the asserted support. The distinguished measure on $-K_Y+P_0$
equals the pushforward of the smooth anticanonical volume almost
everywhere, so its local integral is finite, including at boundary
intersections. This proves both assertions on the flattened graph,
without requiring toroidalization.
\end{proof}

\subsection{A supporting class and finite horizontal support}

\begin{lemma}\label{g:cf:supporting-class}
Assume the data of Section~\ref{g:cf:minima} and suppose no positive power of \(L\) has a nonzero
invariant global section. Then there is a nonzero movable class \(\alpha\)
on \(Y\) satisfying
\[
 \alpha\cdot D=\alpha\cdot D^*=0,
 \qquad \alpha\cdot\Gamma=0
             \quad\text{for every omitted prime }\Gamma.
\]
In particular, \(d^*_\Gamma=d_\Gamma\) whenever
\(\alpha\cdot\Gamma>0\).
\end{lemma}

\begin{proof}
The bounded psh metric makes \(D\) pseudoeffective.  Alternatively,
apply Nadel vanishing and regularity to its divisible multiples:
local boundedness makes their multiplier ideals trivial, and
\(mD+A\) is effective for fixed sufficiently positive \(A\).

The divisor \(D^*\) cannot be rationally linearly equivalent to an
effective rational divisor modulo the span of omitted primes.
Otherwise, after clearing denominators, there would be \(N>0\)
and \(g\in K^*\) with
\[
           N d^*_\Gamma+\operatorname{ord}_\Gamma(g)\geq0
                 \quad\text{at every nonomitted prime}.
\]
Then \(s^Ng\) would have no pole along any prime divisor on \(X\).
For horizontal divisors this follows from the generic regularity of $s$.
For a vertical divisor, its nonexceptional strict transform on
\(W\) dominates a nonomitted prime of \(Y\), and the displayed
inequality and the minimum defining \(d^*_\Gamma\) give the result.
Thus \(s^Ng\) would be a nonzero invariant global section.

Let \(\mathcal S\subset N^1(Y)_{\mathbb R}\) be the span of omitted primes,
and project the cone of effective divisor classes to
\(N^1(Y)_{\mathbb R}/\mathcal S\).  Denote this convex cone by \(C\).
Its closure has nonempty interior, since the ample cone already has
nonempty interior before projection.  The image of \([D^*]\)
belongs to \(\overline C\), by pseudoeffectivity of \(D\) and
effectivity of \(D^*-D\).  It cannot lie in its interior.  Indeed,
the interior of a convex cone agrees with that of its closure and
is contained in the cone itself.  Membership of this rational point
in \(C\) would express it as a finite nonnegative real combination
of rational effective classes.  The corresponding finite rational
linear system then has a nonnegative rational solution.  Since
numerical and rational linear equivalence agree on \(Y\), that
would give the rational effectivity just excluded.

A supporting functional at this boundary point is nonzero,
nonnegative on \(C\), and zero on the image of \([D^*]\).
Pull it back to \(N^1(Y)_{\mathbb R}\).  It is nonnegative on
effective classes and annihilates \(\mathcal S\); by cone duality
\cite{BDPP2013}, it is a nonzero movable class \(\alpha\).
Pseudoeffectivity of \(D\) and effectivity of \(D^*-D\) give
\(0\leq\alpha\cdot D\leq\alpha\cdot D^*=0\).  The effective
difference therefore has zero degree; each of its components has
zero degree.  This proves the final assertion as well.
\end{proof}


For clarity, the numerical-to-linear step in this proof uses
$\operatorname{Pic}^0(Y)=0$: the kernel of the map from the N\'eron--Severi
group to numerical divisor classes is torsion, so rational numerical
relations among the indicated divisors are rational linear relations.
The movable class just constructed annihilates $P_0$ in
Lemma~\ref{g:cf:finite-volume-base}. Apply
Proposition~\ref{g:bt:weighted-slope} with that effective correction, its
integrable distinguished section, and $\alpha\cdot P_0=0$. Every
cotangent line has nonpositive $\alpha$-degree.
Lemma~\ref{g:determinant}, applied to a Cartier multiple of $D$,
gives a fixed cotangent determinant and effective integral divisors
\begin{equation}\label{g:cf:null-divisors}
 N_m\sim M+mD,\qquad m\longrightarrow\infty,\qquad
 \alpha\cdot N_m=0,
\end{equation}
for one fixed line bundle $M$. All indices clear the denominators of $D$.

The field $K_\alpha$ from Lemma~\ref{g:cf:null-field} is the same smaller
field that maximal zero-degree systems give. We now use a different
section-extraction argument on that field, keeping its separate input:
finitely many horizontal null primes. This is the relative counterpart
of the finite-support comparison in \cite[Lemma~2.1]{LP2018}; the
uniqueness verification after extraction is an additional step here.

\begin{proposition}[Finite-support extraction]\label{g:cf:dickson}
Under the data of Section~\ref{g:cf:minima}, assume no positive multiple of
$L$ has an invariant global section. There are a positive integer $a'$ and
an invariant rational section $s'$ of $a'L$ satisfying both the generic
regularity condition and either uniqueness condition of
Lemma~\ref{g:cf:uniqueness} over $K_\alpha$. Moreover
$\operatorname{trdeg}K_\alpha<\operatorname{trdeg}K$.
\end{proposition}

\begin{proof}
Choose one higher equivariant resolution \(\widehat Z\to Z\) on
which the rational map for \(K_\alpha\) is also a morphism.  Pull back
the divisors \(N_m\).  Their parts horizontal for \(K_\alpha\) have
support in a fixed finite set of prime divisors on \(\widehat Z\).
Indeed, a horizontal component mapping onto a prime of \(Y\)
must lie over one of the finitely many horizontal null primes in
Lemma~\ref{g:cf:null-field}.  There are only finitely many components
over those primes.  A prime whose image in \(Y\) has codimension
at least two must be exceptional over \(W\), by the fibre dimension
bound, and there are only finitely many such divisors on this fixed
resolution.

List the horizontal multiplicities of the pulled-back \(N_m\) as
vectors in \(\mathbb Z_{\geq0}^r\), using this fixed finite support.
Order the indices \(m\) increasingly.  Dickson's lemma, in its sequence
form (see \cite[Theorem~1.1(1) and Lemma~3.2]{AichingerAichinger2018}), gives
\(m_2>m_1\) such that the earlier vector is componentwise at most
the later one.  For completeness, the needed sequence form follows
by induction on \(r\): a sequence of nonnegative integers either
has an infinitely repeated value or has a strictly increasing
subsequence; apply this successively to each coordinate.  This
also covers empty horizontal support.  Hence
\(N_{m_2}-N_{m_1}\), pulled back to \(\widehat Z\), is effective
on every \(K_\alpha\)-horizontal divisor.

Put \(b=m_2-m_1>0\).  From \eqref{g:cf:null-divisors}, choose
\(g_0\in K^*\) satisfying
\[
          N_{m_2}-N_{m_1}=bD+\operatorname{div}_Y(g_0).
\]
Set
\begin{equation}\label{g:cf:new-section}
                         a'=ba,\qquad s'=s^bg_0.
\end{equation}
The section \(s'\) is invariant.  On \(\widehat Z\), its divisor
is the pullback of \(b(\operatorname{div}_Z(s)-g_Z^*D)\) plus the pullback of
\(N_{m_2}-N_{m_1}\).  It is therefore nonnegative along all
\(K_\alpha\)-horizontal divisors.  The finitely many remaining poles
lie over proper closed subsets of the new base, so a dense open
avoids them.  This proves generic regularity of the new section.

It remains to verify the ratio requirement, which is essential for
another iteration.  Let \(k\geq0\), and let \(t\) be an invariant
nonzero rational section of \(ka'L\) without poles along prime
divisors on \(X\) horizontal for \(K_\alpha\).  Since \(K_\alpha\subseteq K\),
every \(K\)-horizontal valuation is \(K_\alpha\)-horizontal.  The old
fixed-section uniqueness condition therefore gives
\[
                             g=\frac{t}{s^{kb}}\in K.
\]
Let \(\Gamma\) be a non-null prime of \(Y\).  It is nonomitted,
and \(d^*_\Gamma=d_\Gamma\) by
Lemma~\ref{g:cf:supporting-class}.  Each nonexceptional component
\(Q\) over it on \(W\) is horizontal for \(K_\alpha\): its restriction
to \(K_\alpha\) is \(e_Q\operatorname{ord}_\Gamma=0\).  Its valuation
is the valuation of a prime divisor on \(X\).  Consequently
\[
  e_Q\operatorname{ord}_\Gamma(g)+kb\operatorname{ord}_Q(s)
                         =\operatorname{ord}_Q(t)\geq0.
\]
Taking the minimum over these nonexceptional components proves
\begin{equation}\label{g:cf:ratio-inequality}
                   \operatorname{ord}_\Gamma(g)+kb\,d_\Gamma
                   \geq0.
\end{equation}
The rational divisor \(\operatorname{div}_Y(g)+kbD\) has total
\(\alpha\)-degree zero.  Its coefficients at every non-null prime
are nonnegative by \eqref{g:cf:ratio-inequality}, and the remaining
primes have degree zero.  Each of those nonnegative coefficients
must therefore be zero.  Moreover,
\(\operatorname{div}_Y(g_0)+bD=N_{m_2}-N_{m_1}\) is supported on
null primes.  Subtracting \(k\) times this identity shows that
\[
       \frac{t}{(s')^k}=\frac{g}{g_0^k}\in K_\alpha.
\]
This includes \(k=0\), so fixed-section uniqueness is preserved in
its full stated range.


Strict transcendence-degree decrease follows from
Lemma~\ref{g:cf:null-field}. The regularity just proved holds on the
complete generic fiber and hence survives further modification. This
completes the new data, including the original-prime uniqueness test.
\end{proof}

For the same supporting movable class, affine interpolation in
Lemma~\ref{g:cf:horizontal-slope} and
Proposition~\ref{g:cf:dickson} therefore reach the same field by different
arguments. The first proves that the horizontal multiplicities are affine
functions of the twist. The second only needs finite horizontal support
and chooses a comparable pair of multiplicity vectors. Neither method
replaces the subsequent uniqueness verification.

\subsection{A toric construction of the initial section}

The rational torus quotient can also be initialized using the geometry of
its complete generic fiber. This gives a second construction of existence;
its essential input is toric generation and a weight argument, rather than
the wedge of fundamental vector fields.

\begin{lemma}\label{g:cf:toric-initial-section}
Let an algebraic torus $T$ act on a smooth connected projective variety
$X$, and suppose that $L=-K_X$ is nef. On an equivariant smooth projective
resolution $\pi:Z\to X$, $g:Z\to Y$ of the rational quotient
$\mathbb C(Y)=\mathbb C(X)^T$, some positive power of $\pi^*L$ has a
nonzero invariant section on the complete generic fiber, for its natural
linearization.
\end{lemma}

\begin{proof}
The invariant field is relatively algebraically closed in $\mathbb C(X)$:
an element algebraic over it has a finite orbit, and the divisible abelian
group $T(\mathbb C)$ has no nontrivial finite quotient. Rosenlicht's
quotient theorem gives a dense orbit on the smooth connected general
fiber $F$ of $g$ \cite{Rosenlicht1956}. The stabilizer of that orbit fixes
its closure, so $F$ is toric for an effective quotient of $T$. The nef
bundle $\pi^*L|_F$ has a basepoint-free positive multiple by the toric
basepoint-freeness theorem \cite[Theorem~6.3.12]{CLS2011}.

For any one-parameter subgroup $\lambda$, the Bia\l{}ynicki-Birula
decomposition of $X$ has an open attracting stratum as $t\to0$, whose
fixed component has nonnegative tangent weights, and an open attracting
stratum as $t\to\infty$, whose tangent weights are nonpositive
\cite{BB1973}. Choose $F$ very general so that its image meets both
strata for every $\lambda$; there are only countably many integral
cocharacters. The corresponding orbit limits on the proper $F$ give
fixed points with the two weak signs for the fiber weights of $\pi^*L$,
since $L=\det T_X$ has its natural action. Evaluation of a basepoint-free
multiple at either point is equivariant and surjective. Thus its section
space has weights of both weak signs for every $\lambda$.

It follows that zero is in the convex hull of the torus characters in
this section space: otherwise a strict separating functional could be
chosen rational and then integral. Choose a rational convex combination
with sum zero and clear denominators. The product of the corresponding
nonzero eigensections is a nonzero invariant section of a positive power;
nonzeroness uses integrality of $F$. This is the invariant-sections
construction used in the toric setting of \cite[Section~3]{Muller2025}.

For each positive integer $m$, cohomology and base change and torus
weight splitting identify the invariant section space of $m\pi^*L$ on
fibers over a dense open with a vector bundle of its generic rank.
Choose the very general fiber also in all these countably many opens.
The invariant section just found therefore implies a nonzero invariant
generic rank for its particular exponent. This yields the required
section over $\mathbb C(Y)$. The empty-dimensional fiber and the trivial
torus are included in the argument.
\end{proof}

\begin{remark}
The same weight argument can be made directly on the geometric generic
fiber. Its generic point avoids the complements of the two attracting
strata for every chosen one-parameter subgroup, and properness supplies
the orbit limits. Toric generation and separation of the character
polytope again produce a weight-zero product. Invariants for a split
torus commute with field extension, so this section descends from the
algebraic closure of \(\mathbb C(Y)\). This gives a geometric-generic
form of the preceding construction, without choosing a very general
closed fiber.
\end{remark}

\section{Evaluation metrics and fields with controlled poles}
\label{g:field-variants}

The previous section extracted a generic section from divisors of zero
movable degree. We now vary the choice of the smaller base. Evaluation
in an adjoint direct image supplies a transverse slope estimate.
We apply it to systems generated by four fixed divisor classes,
checking that their maximality still restores uniqueness after
contraction. A separate construction then uses a field of functions
with controlled poles to encode the divisors that may be absorbed
on the original variety.

\subsection{Evaluation metrics and cotangent slopes}
\label{g:cf:evaluation-subsection}

There is a second way to obtain the slope inequality needed for
contraction.  It uses evaluation in an adjoint direct image over a
\emph{further} rational quotient of the base.  This gives a metric on
the determinant of its transverse tangent sheaf, with a correction
supported only on the omitted base primes.

\begin{proposition}\label{g:cf:evaluation-slope}
Let \(X,Y\) be smooth connected projective varieties, let
\(L=-K_X\) have a smooth semipositive metric, and let an algebraic
torus \(T\) act on \(X\).  Suppose that \(f:X\dashrightarrow Y\) is
dominant and invariant and that every invariant rational function on
\(X\) having no poles at the \(f\)-horizontal prime divisors of
\(X\) belongs to \(\mathbb C(Y)\).  Choose a normal projective graph
model
\[
 X\xleftarrow{p}W\xrightarrow{\varphi}Y
\]
on which every vertical prime maps onto a prime divisor of \(Y\).
Let \(N\) be the reduced sum of the primes of \(Y\) above which all
prime divisors of \(W\) are exceptional over \(X\).

For every dominant rational map \(g:Y\dashrightarrow U\) to a smooth
projective variety, let \(\mathcal F\subset T_Y\) be the saturated
sheaf tangent to its general fibers.  There is an integer \(c\geq0\)
such that
\[
 \det(T_Y/\mathcal F)+cN
 \quad\text{is pseudoeffective}.
\]
Consequently, if \(\alpha\ne0\) is a movable curve class on \(Y\)
with \(\alpha\cdot N=0\), then every line subsheaf
\(A\subset(\Omega_Y^1)^{\otimes j}\), \(j\geq0\), satisfies
\(\alpha\cdot c_1(A)\leq0\).
\end{proposition}

The uniqueness assumption in this statement is precisely the
degree-zero instance of Lemma~\ref{g:cf:uniqueness}.  In particular, the
evaluation argument does not require a nef anticanonical class on
\(Y\).

\begin{proof}
Average the metric on \(L\) over the compact form of \(T\), so that
it is invariant.  When \(U\) is a point the determinant assertion is
immediate.  Otherwise replace \(U\) birationally, without changing
\(\mathcal F\), so that every \(g\)-vertical prime \(P\subset Y\)
maps onto a prime \(Q\subset U\).  To arrange this, flatten the
dominating graph and resolve the base \cite{RaynaudGruson1971}.
Each $g$-vertical prime of $Y$ has a strict transform on the
resulting equidimensional source. Properness defines the rational
lift from $Y$ at its generic point, so the resulting map sends
that prime onto a base prime. Put
\(e_P=\operatorname{ord}_P(g^*Q)\).  At the general point of \(P\),
the map has coordinates
\[
 (x,y_2,\ldots,y_{\dim Y})
 \longmapsto (x^{e_P},y_2,\ldots,y_{\dim U}).
\]
The determinant of the differential, after quotienting by its
saturated kernel, therefore has order \(e_P-1\).  Along a horizontal
prime it has order zero.  Thus, in codimension one and hence as line
bundles on \(Y\),
\begin{equation}\label{g:cf:evaluation-determinant}
 \det(T_Y/\mathcal F)=-g^*K_U-R_g,
 \qquad R_g=\sum_{P\text{ vertical}}(e_P-1)P.
\end{equation}
The pullback here is defined on the domain of the rational map and
extended across its codimension-two complement.  The sum is finite,
since its coefficients are the divisorial orders of one determinant
map.

Take a smooth projective equivariant resolution dominating \(W\)
and resolving both rational maps, with morphisms
\[
 \pi:Z\longrightarrow X,\qquad q:Z\longrightarrow Y,
 \qquad r:Z\longrightarrow U.
\]
Write \(e_\pi\) for the invariant Jacobian section of
\(K_Z+\pi^*L\).  On a dense open of \(U\), the sheaf
\[
 \mathcal E=\bigl(r_*(K_{Z/U}+\pi^*L)\bigr)^T
\]
is locally free, commutes with base change, and has its adjoint
\(L^2\) metric.  Berndtsson's theorem gives Griffiths semipositivity
\cite{Berndtsson2009}.  The same is true of the invariant summand:
the weight-zero projection is holomorphic and orthogonal, because
the compact torus preserves the coefficient metric and the fiber
integral.  Its metric is therefore the quotient metric.  This
argument concerns a smooth coefficient metric on the smooth-family
open.

Choose a rational frame for \(\mathcal E\) and a rational frame
\(\eta\) of \(K_U\).  Its sections on \(Z\) have the form
\begin{equation}\label{g:cf:evaluation-coefficients}
 (q^*\phi_i)e_\pi/(r^*\eta),\qquad
 \phi_i\in\mathbb C(Y).
\end{equation}
Indeed, dividing first by the Jacobian and the inverse base frame
gives invariant rational functions on \(X\).  These have no poles
at the primes dominating \(Y\): such primes also dominate \(U\),
and the Jacobian is a unit at their strict transforms.  The assumed
degree-zero uniqueness places the functions in \(\mathbb C(Y)\).

After shrinking dense opens, evaluation of the coefficients in
\eqref{g:cf:evaluation-coefficients} defines a holomorphic functional
on \(g^*\mathcal E\).  For \(v\in\mathcal E_{g(y)}\), denote its
coefficient at \(y\) by \(\phi_v(y)\).  In the frame
\(g^*\eta^{-1}\), set
\begin{equation}\label{g:cf:evaluation-weight}
 \psi(y)=\log\sup_{\|v\|_{g(y)}=1}|\phi_v(y)|^2.
\end{equation}
This is psh, being the log squared norm of a holomorphic section of
the dual of a Griffiths-semipositive bundle.  Replacing \(\eta\) by
\(a\eta\) adds \(\log|g^*a|^2\) to \(\psi\), exactly the
transformation law for a weight on \(-g^*K_U\).  The section
\(e_\pi/(r^*\eta)\) has coefficient \(1\), so this weight is not
identically \(-\infty\).

It remains to account for ramification and for the Jacobian zeros
above the boundary.  Fix a prime \(P\subset Y\), choose a prime
\(D\subset Z\) mapping onto it, and put
\[
 l=\operatorname{ord}_D(q^*P),\qquad
 d_\pi=\operatorname{ord}_D(e_\pi).
\]
If \(P\not\subset N\), choose \(D\) nonexceptional over \(X\),
so that \(d_\pi=0\).  Suppose first that \(P\) maps onto a prime
of \(U\).  With \(j=le_P\), suitable coordinates around a general
point of \(D\) express \(r\) as
\[
 (\tau,z_2,\ldots,z_n)\longmapsto
 (\tau^j,z_2,\ldots,z_h),
 \qquad h=\dim U,\quad D=(\tau=0).
\]
Take \(\eta\) regular and nonvanishing on this base patch.  If a
fiber section \(v\) has coefficient \(a_v\) in a regular frame of
\(K_{Z/U}+\pi^*L\), its adjoint fiber norm contains the density
\[
 \frac{|a_v|^2}{|j\tau^{j-1}|^2}
 |dz_{h+1}\cdots dz_n|^2
\]
up to positive bounded factors with bounded inverse.  The
denominator comes from dividing the total coordinate top form by
the pulled-back base form.  For \(\|v\|=1\), holomorphic submean
on a smaller fiber polydisc gives
\( |a_v|\leq C|\tau|^{j-1}\).
When the fiber dimension is zero, its point contribution to the
norm gives the same bound.  Dividing by the Jacobian coefficient,
which is \(\tau^{d_\pi}\) times a unit, yields
\begin{equation}\label{g:cf:evaluation-boundary}
 |\phi_v(q(z))|\leq C'|\tau(z)|^{j-1-d_\pi}.
\end{equation}
If \(P\) dominates \(U\), then \(r\) is submersive at a general
point of \(D\).  Submean gives instead
\( |\phi_v(q(z))|\leq C'|\tau(z)|^{-d_\pi}\).
In both cases the constants are uniform for unit-norm sections.

Let \(x\) define \(P\) near a general point.  On the chosen patch,
\(q^*x\) is a unit times \(\tau^l\).  Generic submersivity along
\(D\to P\), together with this transverse power, shows that the
patch maps onto a neighborhood of that point of \(P\).  We can
therefore apply the preceding estimates over each nearby \(y\),
with \(|\tau|\) comparable to \(|x(y)|^{1/l}\).  They give
\begin{equation}\label{g:cf:evaluation-exponents}
 \psi\leq O(1)+
 \begin{cases}
 \displaystyle\frac{le_P-1-d_\pi}{l}\log|x|^2,
     &P\text{ vertical over }U,\\[4pt]
 \displaystyle-\frac{d_\pi}{l}\log|x|^2,
     &P\text{ horizontal over }U.
 \end{cases}
\end{equation}
For \(P\not\subset N\), subtracting the weight of \(R_g\) leaves
a locally bounded-above function: in the vertical case the
remaining coefficient is \(1-1/l\geq0\), and in the horizontal
case it is zero.  At a component of \(N\), adding
\(c\log|x|^2\) makes the remaining coefficient nonnegative once
\[
 c\geq (1+d_\pi)/l-1\quad\text{in the vertical case},
 \qquad c\geq d_\pi/l\quad\text{in the horizontal case}.
\]
There are finitely many omitted primes, so one nonnegative integer
\(c\) works for all of them.  The adjusted weights are psh off the
divisors, since the logarithms there are pluriharmonic.  Their upper
bounds extend them across the general points of every boundary
divisor.  The remaining exceptional sets have codimension at least
two; psh removability extends the weights there as well.  The frame
changes identify the resulting metric as one on
\(-g^*K_U-R_g+cN\).  This proves pseudoeffectivity and hence
\begin{equation}\label{g:cf:evaluation-quotient-degree}
 \alpha\cdot c_1(T_Y/\mathcal F)\geq0
 \qquad(\alpha\text{ movable},\ \alpha\cdot N=0).
\end{equation}

We finish with the movable-slope argument of Ou, using the tensor
slope theorem of Greb--Kebekus--Peternell and the algebraicity
criterion of Campana--P\u{a}un
\cite{GKP2016,CampanaPaun2019,Ou2023,LMPTX2023}.
Suppose that a line in \((\Omega_Y^1)^{\otimes j}\) has positive
\(\alpha\)-degree.  For \(j=0\) this is impossible, since a line
subsheaf of \(\mathcal O_Y\) has an effective negative.  For
\(j>0\), tensor slopes and duality imply
\(\mu_{\alpha,\min}(T_Y)<0\).  Applying
\eqref{g:cf:evaluation-quotient-degree} to the identity of \(Y\)
gives \(\alpha\cdot c_1(T_Y)\geq0\).  Thus the positive-slope
part \(\mathcal F_+\) of its Harder--Narasimhan filtration is
nonzero and proper.  Its minimal slope is positive, whereas
\(T_Y/\mathcal F_+\) has nonpositive maximal slope and strictly
negative degree.

The bracket map
\(\mathcal F_+\otimes\mathcal F_+\to T_Y/\mathcal F_+\)
vanishes by the slope gap, so \(\mathcal F_+\) is a foliation.
The positive-minimal-slope criterion makes it algebraically
integrable.  Its rational quotient contradicts
\eqref{g:cf:evaluation-quotient-degree}.  These uses of slope theory
also apply to a boundary movable class.  Equivalently, perturb
toward a rational interior movable class: the strict positive
minimal slope and the bracket slope gap persist.  The tensor
inequality at the original class follows by taking limits from
the interior, using
\[
 \mu_{\alpha+\epsilon\beta,\min}(\mathcal G)
 \geq\mu_{\alpha,\min}(\mathcal G)
       +\epsilon\mu_{\beta,\min}(\mathcal G)
\]
and the corresponding upper bound for maximal slopes.  This
proves the cotangent assertion.
\end{proof}

\subsection{Restricting the contraction to a four-generator semigroup}

The evaluation slope of Proposition~\ref{g:cf:evaluation-slope} can be
used with a smaller collection of linear systems than all null-degree
systems. We record the precise collection because maximality within it
still has to control the poles appearing in the uniqueness test.

\begin{proposition}\label{g:cf:finite-semigroup}
Keep the fixed-section data of Section~\ref{g:cf:minima} and assume
$H^0(Y,\Omega_Y^q)=0$ for $q>0$. After taking a power of $s$, let $D,D^*$
be integral, put $G=D^*-D\geq0$, and let $N$ be the reduced omitted-prime
divisor. Suppose $D$ has a bounded semipositive metric and that every
cotangent line has nonpositive degree against every movable class
annihilating $N$. If there is no invariant global section of a positive
multiple of $L$, a smaller base preserving an invariant generic section
of a positive multiple of $L$ and the original-prime uniqueness condition
is obtained from one of the systems
\[
 |mD+iM+kG+lN|,\qquad m,i,k,l\in\mathbb Z_{\geq0},
\]
where $M$ is a single cotangent determinant line and
$H^0(Y,M+m_jD)\ne0$ for unbounded positive $m_j$.
\end{proposition}

\begin{proof}
The divisor $D+G+N=D^*+N$ cannot be big. A section of a positive
multiple, represented by $\gamma\in\mathbb C(Y)^*$, would satisfy
$\operatorname{ord}_P\gamma+j d_P^*\geq0$ at every nonomitted prime.
The minimum defining $d_P^*$ then makes $s^j\gamma$ regular at every
vertical prime of $X$; its horizontal regularity is already known.
This would be an invariant global section. Since $D$ is pseudoeffective
and $G,N$ are effective, cone duality supplies $\alpha\ne0$ movable with
\[
 \alpha\cdot D=\alpha\cdot G=\alpha\cdot N=0.
\]
The assumed cotangent estimate and the determinant proof in
Lemma~\ref{g:cf:effective-sequence} produce $M$ and effective
$C_j\sim M+m_jD$; effectivity forces $\alpha\cdot M=0$.
Thus every displayed system has degree zero, and none gives a generically
finite map on $Y$.

Choose one nonempty displayed system $|\Lambda|$ of maximal image
dimension, and let $Y'$ be a smooth model of its Stein field. Products
remain in the displayed collection. Exactly the product-map argument in
Lemma~\ref{g:cf:maximal-field} therefore shows that ratios in any one
of these systems lie in $\mathbb C(Y')$. Moreover $\dim Y'<\dim Y$.
The interpolation of Lemma~\ref{g:cf:horizontal-slope} uses only products
of the sections of $M+m_jD$, so it applies within this collection. For
$m_1>m_0$, the ratio of those sections gives a rational section of
$(m_1-m_0)D$ regular on the complete generic fiber over $Y'$. Write $\delta=m_1-m_0$ and let $u$ be the normalized section of
$aL-f^*D$. If $\xi$ is the generic section of $\delta D$ just obtained,
then $u^\delta f^*\xi$ is the new invariant generic section of
$a\delta L$.

We check all-degree original-prime uniqueness directly. Let $w_1,w_2$
be invariant rational sections of $bL$, $b\geq0$, regular at the primes
of $X$ horizontal over $Y'$. Old uniqueness gives
\[
 w_i^a=s^b f^*\gamma_i,\qquad \gamma_i\in\mathbb C(Y)^*.
\]
At a nonomitted prime $P$ of $Y$ horizontal over $Y'$, every
nonexceptional component above it is horizontal over $Y'$ too. Testing
regularity there and taking the nonexceptional minimum gives
\[
 \operatorname{ord}_P\gamma_i+b(d_P+\operatorname{coeff}_P G)\geq0.
\]
Hence, for one sufficiently large $l$, the $\gamma_i$ are rational
sections of $bD+bG+lN$ with poles only at primes vertical over $Y'$.

These finitely many poles can be cleared by the same section of some
$n\Lambda$; if there are no poles, use $n=0$ and the constant section. Resolve $|\Lambda|$. Each such prime maps into a proper
closed subset of its image: a valuation trivial on the image field is
trivial on its algebraic closure, so verticality persists for the image
field. Choose a projective hypersurface containing those images but not
the whole system image. Substituting the defining sections of $\Lambda$
gives a nonzero section vanishing at each pole prime; its sufficiently
high power gives the required orders. The fixed divisor of the resolved
system only adds vanishing. Multiplication by this common section places
both $\gamma_i$ in one system of the displayed semigroup. Maximality
puts their ratio in $\mathbb C(Y')$. Thus
$(w_1/w_2)^a\in\mathbb C(Y')$, and relative algebraic closedness inside
$\mathbb C(X)$ gives $w_1/w_2\in\mathbb C(Y')$.
Lemma~\ref{g:cf:uniqueness} now restores fixed-section uniqueness,
including degree zero. This proves the proposition.
\end{proof}

For the geometric data constructed from a rational quotient, the bounded
metric is supplied by Lemma~\ref{g:cf:finite-volume-base}, while
Proposition~\ref{g:cf:evaluation-slope} supplies exactly the cotangent
hypothesis. This verifies a complete contraction using the evaluation
metric, independently of the signed-volume slope argument. The four-generator
semigroup need not generate the whole null-prime field; its own
maximality and pole-clearing property are what the proof uses.

\subsection{A parameter field allowing poles at omitted primes}
\label{g:cf:parameter-subsection}

The next construction uses the signed-boundary slope inequality
of Proposition~\ref{g:bt:signed-slope}.  Its field is built directly
from rational functions with prescribed orders away from the
omitted primes.  Allowing those poles is useful because no divisor
of the original variety lies above an omitted prime.  The
effective divisors used for horizontal interpolation, however,
remain effective on the whole base.

Use fixed-section data \((f:X\dashrightarrow Y,k,s)\) satisfying
the regularity and uniqueness conditions of
Section~\ref{g:cf:conditions}, and a normal equidimensional model
\(X\xleftarrow pW\xrightarrow\varphi Y\).  For a prime
\(\Gamma\subset Y\), and primes \(A\subset W\) mapping onto it,
write
\[
 e_A=\operatorname{ord}_A(\varphi^*\Gamma),\qquad
 l_A=\operatorname{ord}_A(p^*s),\qquad
 b_\Gamma=\min_{A\to\Gamma}\frac{l_A}{e_A}.
\]
Let \(C_1,\ldots,C_v\) be the omitted primes, above which every
such \(A\) is exceptional over \(X\).  Off this finite list put
\[
 b^\sharp_\Gamma=
 \min_{\substack{A\to\Gamma\\A\text{ nonexceptional over }X}}
 \frac{l_A}{e_A},
\]
and give \(b^\sharp_\Gamma\) the value zero on the omitted primes.
Denote the resulting divisors by \(B,B^\sharp\).  Replace \(s\)
by a power, and scale \(k,B,B^\sharp\), so that both divisors are
integral.  We have
\begin{equation}\label{g:cf:parameter-model}
 D_s:=\operatorname{div}(p^*s)-\varphi^*B\geq0,
 \qquad b^\sharp_\Gamma\geq b_\Gamma
 \quad(\Gamma\ne C_i).
\end{equation}
Moreover, for \(l>0\) integral,
\begin{equation}\label{g:cf:parameter-lift}
 \operatorname{ord}_\Gamma v+l b^\sharp_\Gamma\geq0
 \quad(\Gamma\ne C_i)
 \quad\Longrightarrow\quad
 s^l f^*v\in H^0(X,klL)^T\setminus\{0\}
\end{equation}
for \(v\in\mathbb C(Y)^*\).  Indeed the inequalities test every
vertical prime of \(X\) through its nonexceptional transform on
\(W\); the original horizontal primes have nonnegative orders
by the fixed-section condition.  Normality of \(X\) then gives
regularity.

For the geometric application, assume again that $L=-K_X$ is smoothly
semipositive and $H^0(X,\Omega_X^q)=0$ for $q>0$. Pullback of forms
gives the same vanishing on $Y$, hence $\chi(\mathcal O_Y)=1$.
Prepare this graph toroidally as in Section~\ref{g:cf:preparation};
then define the minima on that prepared graph as above. The bounded metric and signed nef anticanonical class in
Lemma~\ref{g:cf:metric-conclusions}, the determinant construction in Lemma~\ref{g:determinant}, and
Proposition~\ref{g:bt:signed-slope} give an integral divisor \(R\)
and an unbounded set \(\mathcal M\subset\mathbb Z_{>0}\) with
\begin{equation}\label{g:cf:parameter-effective}
 [R]\in\overline{\operatorname{Psef}(Y)
                     +\sum_{i=1}^v\mathbb R[C_i]},
 \qquad N_m=mB-R+\operatorname{div}(q_m)\geq0
 \quad(m\in\mathcal M).
\end{equation}
Here \(q_m\in\mathbb C(Y)^*\), and the effectivity of \(N_m\)
includes the omitted primes. To make the sign of \(R\) explicit,
Lemma~\ref{g:determinant} gives a line \(M\) in a cotangent tensor
power of order \(h\) and effective divisors
\[
 N_m\sim mB+M,\qquad R=-M.
\]
The signed boundary \(B_{\mathrm{sgn}}\) furnished by
Lemma~\ref{g:cf:metric-conclusions} has negative part supported
on the \(C_i\). Proposition~\ref{g:bt:signed-slope} gives
\[
 [R]+h[B_{\mathrm{sgn}}^-]\in\operatorname{Psef}(Y),
\]
so \([R]\) lies in
\(\operatorname{Psef}(Y)+\sum_i\mathbb R[C_i]\), and in particular
in its closure as required in \eqref{g:cf:parameter-effective}.
The following argument depends only on
the explicit data \eqref{g:cf:parameter-model}--
\eqref{g:cf:parameter-effective} and fixed-section uniqueness.

\begin{proposition}\label{g:cf:parameter-contraction}
Assume fixed-section data $(f:X\dashrightarrow Y,k,s)$ as in
Section~\ref{g:cf:conditions}. Let $B,B^\sharp$ be the integral
vertical-minimum divisors of the preceding normal equidimensional graph,
and let $C_i$ be its omitted primes. Assume that an integral divisor $R$
and an unbounded sequence of effective divisors satisfy
\eqref{g:cf:parameter-effective}.
Then either a positive multiple of \(L\) has a
nonzero invariant section on \(X\), or there is a dominant rational
map \(Y\dashrightarrow Y'\), with \(\dim Y'<\dim Y\), such that
its composite with \(f\) admits fixed-section data satisfying
the regularity and uniqueness conditions of
Section~\ref{g:cf:conditions}.  Its base field is relatively
algebraically closed in \(\mathbb C(X)\).
\end{proposition}

\begin{proof}
Fix \(a\in\mathcal M\).  For \(t\in\mathbb Q\), \(t\geq a\),
and \(h\in\mathbb Z_{>0}\), define a vector space of rational
functions by
\begin{equation}\label{g:cf:parameter-spaces}
 V_{t,h}=\{0\}\cup
 \left\{v\in\mathbb C(Y)^*:
 \operatorname{div}(v)+h(tB^\sharp-R)\geq0
 \text{ away from the }C_i\right\}.
\end{equation}
Rational coefficients cause no ambiguity: this is an inequality
of rational divisors.  These spaces need not be finite-dimensional.
Their multiplication law is
\begin{equation}\label{g:cf:parameter-product}
 V_{t_1,h_1}V_{t_2,h_2}
 \subset
 V_{(h_1t_1+h_2t_2)/(h_1+h_2),\,h_1+h_2}.
\end{equation}
By \eqref{g:cf:parameter-effective} and
\(B^\sharp\geq B\) away from the omitted primes,
\(q_m\in V_{m,1}\).  Every rational \(t\geq a\) is a rational
convex combination of \(a\) and some \(m\in\mathcal M\) with
\(m\geq t\).  Products of powers of \(q_a,q_m\) therefore show
that \(V_{t,h}\ne0\) for some \(h>0\).

At a fixed \(t\), let \(F_t\) be the set of ratios \(v/w\) with
\(v,w\in V_{t,h}\), \(w\ne0\), allowing all positive \(h\).
It is a field: multiplication and cross-multiplication for addition
use \eqref{g:cf:parameter-product}, and inverses interchange
numerator and denominator.  Its relative algebraic closure in
\(\mathbb C(Y)\) is independent of \(t>a\).  To see this, take
\(a<t_1<t_2\) and a ratio at \(t_1\).  Pad equal powers of its
numerator and denominator with a common power of \(q_m\), where
\(m>t_2\).  Choose the powers so that their weighted-average
parameter is \(t_2\).  A positive power of the original ratio
then belongs to \(F_{t_2}\).  Conversely, pad a ratio at \(t_2\)
with \(q_a\) to reach \(t_1\).  Clearing the rational weights
makes all the powers integral.  Each field is therefore contained
in the relative algebraic closure of the other.

Write \(K\) for the common relatively closed field, and choose a
smooth projective model \(Y'\) with \(\mathbb C(Y')=K\).  This
field is finitely generated: choose a transcendence basis of
\(K/\mathbb C\); the algebraic closure of the resulting rational
function field inside the finitely generated field \(\mathbb C(Y)\)
is finite over it.  Since \(\mathbb C(Y)\) is relatively
algebraically closed in \(\mathbb C(X)\), the same is true of
\(K\).  We now distinguish whether this field has lost dimension.

\emph{The full-dimensional case.}
Suppose \(\operatorname{trdeg}_{\mathbb C}K=d=\dim Y\).
At one fixed rational \(t>a\), finitely many ratios in \(F_t\)
are algebraically independent of cardinality \(d\).  Use common
products and powers to place their numerator and denominator
sections in one degree.  A sufficiently large rational divisor
\(C\) supported on the \(C_i\) compensates their finitely many
omitted-prime poles.  After clearing denominators, these sections
belong to a multiple of \(tB^\sharp-R+C\) and give a rational map
with \(d\)-dimensional image.  Hence
\begin{equation}\label{g:cf:parameter-big}
 tB^\sharp-R+C\quad\text{is big}.
\end{equation}

The closure in the first condition of
\eqref{g:cf:parameter-effective} is sufficient to remove \(R\)
from this assertion.  Choose pseudoeffective classes \(P_j\) and
real divisors \(E_j\) supported on the \(C_i\) such that
\(P_j+[E_j]\to[R]\).  Then
\[
 [tB^\sharp+C-E_j]-P_j
 \longrightarrow [tB^\sharp-R+C].
\]
The classes on the left are big for large \(j\), by openness of
the big cone.  Adding \(P_j\) preserves bigness, so
\(tB^\sharp+C-E_j\) is big.  Rational approximation of the
coefficients on the finite support yields a rational divisor
\(C'\), supported on that same list, with \(tB^\sharp+C'\) big.
Consequently, for some positive integer \(h\) clearing all
denominators and some \(v\in\mathbb C(Y)^*\),
\[
 \operatorname{div}(v)+h(tB^\sharp+C')\geq0.
\]
Away from the omitted primes this is precisely the hypothesis of
\eqref{g:cf:parameter-lift}, with the positive integer \(l=ht\).
It produces the asserted global invariant section.

\emph{The smaller-base case.}
Suppose now that \(\dim Y'<d\).  Choose \(b\in\mathcal M\) with
\(b>a\), write \(\delta=b-a\), and set
\begin{equation}\label{g:cf:parameter-new-section}
 f':X\dashrightarrow Y',\qquad
 k'=k\delta,\qquad
 s'=s^\delta f^*(q_b/q_a),
\end{equation}
where \(f'\) is the composite rational map.  Both the map and
section are invariant.  We verify separately that the section is
regular over the generic point and that it restores uniqueness.

For every \(m\in\mathcal M\) with \(m>b\), the numerator and
denominator in
\begin{equation}\label{g:cf:parameter-interpolation}
 \frac{q_m^{\delta}q_a^{m-b}}{q_b^{m-a}}
\end{equation}
belong to \(V_{b,m-a}\), since
\(\delta m+(m-b)a=b(m-a)\).  Their ratio therefore lies in
\(K\).  On a common smooth resolution mapping to \(W,Y,Y'\),
let \(n_m\) be the coefficient of the pullback of \(N_m\) at a
prime horizontal over \(Y'\).  Its valuation is trivial on
\(K\), so \eqref{g:cf:parameter-interpolation} implies
\[
 \delta(n_m-n_a)=(m-a)(n_b-n_a).
\]
Every \(n_m\geq0\), and the indices \(m\) are unbounded.
It follows that \(n_b-n_a\geq0\).  The divisor of the pullback
of \(s'\) is
\[
 \delta D_s+\varphi^*(N_b-N_a),
\]
with pullbacks to the common resolution understood.  Its
horizontal coefficients are nonnegative by
\eqref{g:cf:parameter-model} and the preceding inequality.
Thus \(s'\) has no poles over the generic point of \(Y'\),
including at exceptional horizontal primes.

To prove uniqueness, take \(r\geq0\) and a nonzero invariant
rational section \(\tau\) of \(rk'L\) having no poles on the
prime divisors of \(X\) horizontal over \(Y'\).  These include
the primes horizontal over \(Y\).  The old fixed-section
condition writes
\[
 \tau=s^{\delta r}f^*u,\qquad
 (s')^r=s^{\delta r}f^*u_s,
 \qquad u_s=(q_b/q_a)^r,
\]
with \(u,u_s\in\mathbb C(Y)^*\).  At every non-omitted prime
\(\Gamma\) whose valuation is trivial on \(K\), the
nonexceptional primes above it are horizontal over \(Y'\).
Regularity of both sections there gives
\begin{equation}\label{g:cf:parameter-horizontal-tests}
 \operatorname{ord}_\Gamma(u)+\delta r b^\sharp_\Gamma\geq0,
 \qquad
 \operatorname{ord}_\Gamma(u_s)+\delta r b^\sharp_\Gamma\geq0.
\end{equation}
Only finitely many non-omitted primes fail either inequality.
Every such prime has valuation nontrivial on \(K\).

We can compensate all these deficits with a common element of
one \(V_{b,h}\).  For a non-omitted prime \(\Gamma\), define
the excess of \(v\in V_{b,h}\setminus\{0\}\) by
\[
 \epsilon_\Gamma(v)=
 \operatorname{ord}_\Gamma(v)
   +h(bb^\sharp_\Gamma-\operatorname{coeff}_\Gamma R)
 \geq0.
\]
If \(\operatorname{ord}_\Gamma\) is nontrivial on \(K\), it is
nontrivial on \(F_b\): a valuation trivial on a field remains
trivial on an algebraic extension.  Some ratio \(v_1/v_2\), with
\(v_1,v_2\in V_{b,h_\Gamma}\), therefore has nonzero order.
The two nonnegative excesses are different, so at least one is
strictly positive.  Choose that element.  Products of sufficiently
large powers of these chosen elements, one for each deficit
prime, give \(v\in V_{b,h}\setminus\{0\}\) whose excess
compensates every deficit of both \(u\) and \(u_s\).
At all other non-omitted primes, the excesses remain nonnegative.
If there are no deficits, any nonzero element of \(V_{b,h}\)
suffices.  Thus
\[
 uv,\ u_s v\in V_{\,b+\delta r/h,\,h}.
\]
The parameter is still greater than \(a\), also when \(r=0\).
It follows that
\[
 \tau/(s')^r=u/u_s\in F_{b+\delta r/h}\subset K.
\]
This is the required fixed-section uniqueness for the smaller
base.

These verifications also retain the adjoint rank conditions:
on a smooth resolution, multiplying powers of \(s'\) by the
Jacobian gives a nonzero invariant adjoint section in every
nonnegative multiple of $k'$.  Dividing any other such section by the
Jacobian gives a rational section regular at the original
horizontal primes; the uniqueness just proved makes it a
multiple over \(K\) of that generator.  Thus the contraction
returns precisely the section and rank data needed for the next
step, including degree zero.
\end{proof}
